\scanpage{26}

en associant à $(g,(\alpha_i)_{i\in\{0,1,\infty\}})$ le système \deleted{des} $p=1$, $(g_i)=\deleted{\unsure{(l_i^{\alpha_i})}}\deleted{\unsure{(l_i^{\alpha_i}g^{-1})}}(g\,l_i^{\alpha_i})$

% slanted
\margin{\uline{NB} Il y a lieu d'introduire un groupe $S\widehat{\unsure{\Sigma|E}}$ extension de $\Sigma$ par $\mathbb Z^{\{0,1,\infty\}}$ (déduite d'une extension canonique de $\Gamma$ par $\mathbb Z^{\{0,1,\infty\}}$) et \unsure{d'épingler|d'expliciter|d'introduire} les relèvements des \deleted{\unsure{de}} $\sigma_i,\rho$ etc en des éléments \inserted{déjà} $\widetilde\sigma_i,\widetilde\rho$ de $S\Sigma$ \ldots.}

Il faudrait expliciter les opérations des éléments $\sigma_i$ ou $\varepsilon_i$ \deleted{sur} $(\in S\widehat{\widehat\Sigma}^{!})$ sur $\deleted{\unsure{\Sigma}}\widetilde u\in S\widehat{\widehat\Sigma}^{!}$ \deleted{\unsure{de base}} par conjugaison, et en particulier l'action de $\rho$. On trouve que \deleted{$\chi$ \unsure{(bien sûr)}} les $\gamma_i$ sont ``covariants'' par $\rho$, i.e.

\[
\gamma_{\unsure{i|i'}}(\rho\widetilde u\rho^{-1})=\gamma_{i+1}\unsure{(\bar u)|(\widetilde u)}
\]

\deleted{i.e. si $\widetilde u$ est défini par}
% archive sheet number: 25
